% =====================================================================
%  CV — Apple Machine Learning Research (MLR), Paris
%        Research Internship, ~6 months, onsite Paris, start Feb 2027
%  Dérivé de master-cv.md (source de vérité unique)
%
%  CIBLAGE — labo de ML FONDAMENTAL (Ablin : différentiation implicite,
%  optimisation ; lignée Cuturi : transport optimal ; publie NeurIPS/ICLR).
%  Ce qui remonte, honnêtement :
%   - Le moteur physique différentiable = une COUCHE D'OPTIMISATION SOUS
%     CONTRAINTE différentiable (RNEA, matrice de masse, PD implicite,
%     solveur Gauss-Newton, projection sur le cône) dans laquelle on
%     rétro-propage. C'est le point de contact réel avec le groupe.
%   - Maths from-scratch et algèbre linéaire numérique (SVD, QR, gradient
%     conjugué, Cholesky) : leur socle.
%   - Rigueur expérimentale : le résultat NÉGATIF d'InterMask, publié tel
%     quel, est un signal fort pour un labo de recherche.
%   - Publications en tête (CVPR première autrice).
%
%  CE QUI DESCEND / SORT : IEC 61508, K3s, AWS, MLOps, métriques motion 3D
%  détaillées, explicabilité en section propre. Un labo de ML fondamental
%  ne recrute pas là-dessus.
%
%  CE QUI N'EST PAS REVENDIQUÉ, faute de preuve : transport optimal, théorie
%  de l'optimisation avancée, diffusion. La proximité (optimisation
%  différentiable, couche implicite) est nommée ; la possession, non.
%
%  INTERDITS (master-cv.md §11) : projet clical, −58 %, PFA 82,3 %,
%  audience Medium, Google Final Round, rejet antérieur du papier ICLR,
%  mandarin « working knowledge » (niveau réel : débutant).
%
%  ⚠ AVANT ENVOI : couple Sector 100 %/70 % [À CONFIRMER] §4.3 ; titre exact
%  du papier de Yilin Wang inconnu ; postuler TÔT (l'annonce le dit).
% =====================================================================

\documentclass[10pt, letterpaper]{article}

\usepackage[
    ignoreheadfoot,
    top=0.9 cm,
    bottom=0.8 cm,
    left=1.2 cm,
    right=1.2 cm,
    footskip=0.9 cm,
]{geometry}
\usepackage{titlesec}
\usepackage{enumitem}
\usepackage{fontawesome5}
\usepackage[dvipsnames]{xcolor}
\usepackage{textcomp}

\definecolor{MyRed}{RGB}{160, 20, 40}

\usepackage[
    pdftitle={Melissa Colin — CV — Apple MLR Paris Research Internship},
    pdfauthor={Melissa Colin},
    colorlinks=true,
    urlcolor=MyRed
]{hyperref}

\usepackage{charter}

\pagestyle{empty}
\setlength{\parindent}{0pt}

\titleformat{\section}
  {\large\bfseries\scshape\raggedright\color{MyRed}}{}{0em}{}[\titlerule]
\titlespacing{\section}{0pt}{6pt}{3pt}

\setlist[itemize]{leftmargin=*, noitemsep, topsep=1pt, partopsep=0pt, parsep=0pt}

\newcommand{\entry}[4]{%
  \noindent \textbf{#1} \hfill #2 \\
  \textit{#3} \hfill \textit{#4} \par
}

\IfFileExists{personal.tex}{\input{personal.tex}}{}
\providecommand{\myphone}{+33 7 82 52 XX XX}

\begin{document}

% ---------------------------------------------------------------- HEADER
\begin{center}
    {\fontsize{20}{20}\selectfont \textbf{Mélissa COLIN}} \\[3pt]
    {\large Machine Learning $\cdot$ Differentiable Optimization \& Generative Models} \\[4pt]

    \small{
    \faMapMarker* Bordeaux / Paris, France \quad|\quad
    \faPhone* \myphone{} \quad|\quad
    \faEnvelope~\href{mailto:contact@melissacolin.ai}{contact@melissacolin.ai} \\
    \faGlobe~\href{https://melissacolin.ai}{melissacolin.ai} \quad|\quad
    \faLinkedin~\href{https://www.linkedin.com/in/melissacolin/}{LinkedIn} \quad|\quad
    \faGithub~\href{https://github.com/melissa-colin}{github.com/melissa-colin} \quad|\quad
    \faGraduationCap~\href{https://scholar.google.com/citations?user=7r7iFpsAAAAJ}{Scholar}
    }
\end{center}

\vspace{2pt}

{\small \textit{Final-year MEng student in computer science, with a first-author paper in preparation for CVPR 2027. My work centres on \textbf{differentiable simulation}: I wrote a rigid-body
dynamics engine from scratch and backpropagate through the unrolled simulation to train
neural models, and I care about rigorous experimentation, including publishing negative results. Available
for a 5--6 month research internship, onsite in Paris, starting February 2027.}}

% ---------------------------------------------------------------- PUBLICATIONS
\section{Publications}
\begin{itemize}
    \item \textbf{Colin, M.}, Wang, Y., Cheng, L. \textit{InterForce: Physically Admissible Multi-Person Motion Refinement.} First author, in preparation, targeting \textbf{CVPR 2027}.
    \item \textbf{Colin, M.}, Chraibi Kaadoud, I. \textit{Performances and Explainability of ViT and CNN Architectures: An Empirical Study Using LIME, SHAP and Grad-CAM.} RJCIA / PFIA 2024, La Rochelle. First author, published, presented orally. \hfill \href{https://hal.science/hal-04641791}{[HAL]}
    \item Co-author, \textit{X-Dancers} (3D group choreography), submitted to ICLR 2027.
\end{itemize}

% ---------------------------------------------------------------- RESEARCH
\section{Research Experience}

\entry{Vision and Learning Lab, University of Alberta (Prof.\ Li Cheng)}{May 2026 -- Sep 2026}%
{Visiting Research Student, \textbf{Mitacs Globalink Research Award}}{Edmonton, Canada}
\begin{itemize}
    \item Wrote a \textbf{differentiable rigid-body dynamics engine from scratch} --- recursive Newton--Euler, composite-rigid-body mass matrix, implicit (proximal) integration --- and \textbf{backpropagate through the unrolled simulation}, using \textbf{truncated BPTT} to keep the memory cost of unrolling tractable. Verified against nine analytic invariants to machine precision. (A separate constrained projection onto the friction cone is precomputed offline.)
    \item Cast physical validity as a \textbf{representation constraint rather than a loss penalty}: the model outputs forces confined to the admissible set, so impossible motion is unrepresentable by construction. Share of trajectories a physics simulator reproduces without divergence: \textbf{19\% $\to$ 79\%}.
    \item Designed the accompanying \textbf{permutation-equivariant transformer} (attention factorised over set elements, joints and time), and a \textbf{27{,}000-clip corpus} generated by render-then-estimate.
\end{itemize}

\vspace{3pt}

\entry{ENSEIRB-MATMECA --- Research Initiation elective, 17/20}{Feb -- May 2026}%
{Reproduction and critical study of \textit{InterMask} (ICLR 2025)}{Bordeaux, France}
\begin{itemize}
    \item Reproduced the \textbf{VQ-VAE} generative baseline (FID 1.034 $\to$ \textbf{0.918}), then tested two hypotheses on the constraints in the loss.
    \item \textbf{Reported a negative result:} neither variant improved FID or R-Precision, and I published the finding as-is rather than dropping it.
\end{itemize}

% ---------------------------------------------------------------- SKILLS
\section{Technical Skills}
\begin{itemize}
    \item \textbf{Maths for ML:} \textbf{differentiable simulation}, backpropagation through an unrolled physics rollout (truncated BPTT) $\cdot$ numerical linear algebra (SVD, QR, Cholesky, conjugate gradient) $\cdot$ Gauss--Newton, Newton--Raphson $\cdot$ convex \& integer programming $\cdot$ probability, optimisation, complexity theory.
    \item \textbf{Deep learning:} \textbf{PyTorch} (advanced) $\cdot$ Transformers $\cdot$ \textbf{VQ-VAE} $\cdot$ CNNs $\cdot$ LoRA/PEFT $\cdot$ implemented from scratch: backpropagation, PCA, K-Means, logistic regression, naive Bayes.
    \item \textbf{Programming \& compute:} \textbf{Python} $\cdot$ C $\cdot$ C++ $\cdot$ Java $\cdot$ SQL $\cdot$ Bash $\cdot$ Linux/POSIX $\cdot$ Git $\cdot$ multi-GPU / SLURM-style clusters $\cdot$ \LaTeX{}.
    \item \textbf{Languages:} French (native) $\cdot$ English \textbf{B2, TOEIC 840/990}, research placement conducted entirely in English.
\end{itemize}

% ---------------------------------------------------------------- EDUCATION
\section{Education}

\entry{ENSEIRB-MATMECA, Bordeaux INP}{Sep 2024 -- Expected Sep 2027}%
{MEng in Computer Science (\textit{Diplôme d'ingénieur}), AI track, final year}{Bordeaux, France}
\begin{itemize}
    \item Ranked \textbf{5th/81 (top 6\%)} in semester 8, with \textbf{4th/81 on the project unit} (French /20 scale, transcript available). \textbf{Deep Learning 18/20} $\cdot$ \textbf{Research Initiation 17/20} $\cdot$ numerical algorithms, operations research, complexity theory.
\end{itemize}

\vspace{3pt}

\entry{EPSI Bordeaux}{Sep 2021 -- Sep 2024}%
{BSc in Artificial Intelligence \& Data Science, ranked \textbf{1st/10} in final year}{Bordeaux, France}

% ---------------------------------------------------------------- OTHER
\section{Selected Experience \& Awards}
\begin{itemize}
    \item \textbf{Mitacs Globalink Research Award (2026)} --- competitive international research funding, awarded on project proposal.
    \item \textbf{R\&D Engineer, AI \& Safety}, Sector Group (Jun--Aug 2025): tuned a retrieval pipeline for an abstention regime --- \textbf{100\% precision on the answers returned}, at 70\% recall over 100+ queries.
    \item \textbf{Applied AI Developer}, Cali Intelligences (2023--2024, apprenticeship): built the full training pipeline for detection models, \textbf{accuracy 60\% $\to$ 90\%}.
    \item \textbf{Vice-President, Forum INGENIB}: \texteuro{}100k budget, 20-student team, 800 students and 80 companies hosted.
\end{itemize}

\end{document}
